\section*{Chapitre \ref{ch:b3:cancer-biology} --- Biologie du cancer}

\begin{solution}{exo:b3:cancer-biology:1}
Un \omterm{def:b3:cancer-biology:genes}{oncogène} est une version à gain de fonction d'un gène qui favorise
la prolifération ou la survie~; un seul allèle mutant suffit
(dominant)~: \emph{RAS} (mutation ponctuelle), \emph{MYC}
(translocation, amplification), \emph{HER2}, \emph{BCR--ABL}. Un
suppresseur de \omterm{def:b3:cancer-biology:hallmarks}{tumeur} freine la prolifération ou impose l'arrêt, la
mort ou la réparation~; les deux allèles doivent être perdus (récessif
à l'échelle de la cellule)~: \emph{RB}, \emph{TP53}, \emph{APC},
\emph{BRCA1}.
\end{solution}

\begin{solution}{exo:b3:cancer-biology:2}
Prolifération soutenue~: \emph{RAS}, \emph{MYC}. Insensibilité aux
suppresseurs de croissance~: \emph{RB}, \emph{CDKN2A} (p16). Résistance
à la mort~: \emph{BCL2}, \emph{TP53}. Immortalité réplicative~: le
promoteur de la \omterm{def:b3:cancer-biology:telomerase}{télomérase}. Angiogenèse~: \emph{VEGF} induit par HIF
(perte de \emph{VHL}). \omterm{def:b3:cancer-biology:metastasis}{Invasion} et \omterm{def:b3:cancer-biology:hallmarks}{métastase}~: perte de la
E-cadhérine, les facteurs épithélio-mésenchymateux. Instabilité du
génome~: gènes de la \omterm{prop:b3:dna-repair:fidelity}{réparation des mésappariements}, \emph{\omterm{def:b3:dna-repair:dsb}{BRCA}}.
Métabolisme dérégulé~: \emph{MYC}, \omterm{prop:b3:cancer-biology:pathways}{voie PI3K}. \omterm{prop:b3:cancer-biology:immune}{Échappement immunitaire}~:
\emph{PD-L1}, perte du CMH. Inflammation pro-tumorale~: la
signalisation \emph{NF-$\kappa$B}.
\end{solution}

\begin{solution}{exo:b3:cancer-biology:3}
Les enfants ayant des antécédents familiaux développaient plusieurs
\omterm{def:b3:cancer-biology:hallmarks}{tumeurs} aux deux yeux, et tôt~; les cas sporadiques une seule \omterm{def:b3:cancer-biology:hallmarks}{tumeur},
plus tard. Il en a inféré que deux événements étaient nécessaires dans
une cellule~: les patients familiaux en avaient hérité un et n'avaient
besoin que d'un second, somatique (assez fréquent pour survenir dans
plusieurs cellules rétiniennes), tandis que les patients sporadiques
avaient besoin que les deux surviennent dans une même cellule ---
rare, tardif, unique. D'où un gène dont les deux copies doivent être
perdues~: le premier suppresseur de \omterm{def:b3:cancer-biology:hallmarks}{tumeur}.
\end{solution}

\begin{solution}{exo:b3:cancer-biology:4}
Un mutant de \emph{Salmonella} exigeant de l'histidine est étalé sans
histidine avec le composé à tester~; des colonies n'apparaissent que
sur des cellules où une mutation nouvelle a corrigé le défaut, de sorte
que leur nombre mesure la mutagénicité. On ajoute un extrait de foie
parce que bien des cancérogènes sont inoffensifs tant que des enzymes
hépatiques ne les convertissent pas en métabolites réactifs, comme cela
se produit dans le corps~; les bactéries n'ont pas ces enzymes.
\end{solution}

\begin{solution}{exo:b3:cancer-biology:5}
L'incidence croît d'un facteur $100$ pendant que l'âge double~:
$2^{k-1} = 100$, $k - 1 =
\log_{2}100 = 6.6$, soit environ sept ou huit événements limitants dans
la lecture naïve --- moins si l'on admet une expansion clonale entre
les coups.
\end{solution}

\begin{solution}{exo:b3:cancer-biology:6}
$L = \sqrt{2Dc_{0}/q}$~: pour $q = 0.01$, $\sqrt{2\times 2\times 10^{-9}
\times 0.05/0.01} = \qty{141}{\micro m}$~; pour $q = 0.03$,
\qty{82}{\micro m}. Une \omterm{def:b3:cancer-biology:hallmarks}{tumeur} à croissance rapide respire plus vite,
épuise le dioxygène sur une distance plus courte, et son centre meurt à
une taille moindre.
\end{solution}

\begin{solution}{exo:b3:cancer-biology:7}
$10^{11}\times 10^{-3n} < 1$ exige $n > 3.7$~: quatre cures. Après deux
cures il reste $10^{5}$ cellules --- un dixième de millimètre cube,
bien en deçà des $10^{9}$ détectables --- de sorte que la \omterm{def:b3:cancer-biology:hallmarks}{tumeur} a
«~disparu~» pour tous les examens alors que cent mille cellules
subsistent~; le traitement doit se poursuivre pour le nombre de cures
prévu, et non jusqu'à la disparition de la \omterm{def:b3:cancer-biology:hallmarks}{tumeur}.
\end{solution}

\begin{solution}{exo:b3:cancer-biology:8}
Les deux lésions activent la même voie (récepteur $\to$ \omterm{def:b3:cancer-biology:genes}{Ras} $\to$
ERK)~; une fois l'une présente, l'autre ne confère plus d'avantage et
n'est pas sélectionnée. Un inhibiteur d'EGFR bloque le récepteur~; une
cellule qui acquiert une mutation de KRAS en aval rétablit le signal
alors que le récepteur reste bloqué, et elle est sélectionnée pendant
le traitement --- une résistance par contournement.
\end{solution}

\begin{solution}{exo:b3:cancer-biology:9}
E7 libère E2F de Rb, ce qui fait franchir à la cellule épithéliale
infectée le \omterm{def:b3:cell-cycle-apoptosis:checkpoints}{point de restriction} sans facteurs de croissance~; E6
détruit \omterm{def:b3:cancer-biology:genes}{p53}, de sorte que la prolifération indue et les dégâts qu'elle
cause ne déclenchent ni arrêt ni \omterm{def:b3:cell-cycle-apoptosis:apoptosis}{apoptose}. Deux caractères sont fournis
d'un coup, mais les autres --- immortalité, \omterm{def:b3:cancer-biology:metastasis}{invasion}, instabilité ---
exigent d'autres mutations somatiques, qui mettent des décennies à
s'accumuler dans les cellules infectées de façon persistante. Le vaccin
suscite des anticorps neutralisants contre la capside et empêche
l'infection~; administré avant l'exposition, il n'y a jamais de cellule
infectée à transformer.
\end{solution}

\begin{solution}{exo:b3:cancer-biology:10}
Après le $i$-ième coup le clone croît d'un facteur $g$, de sorte que
$g^{i}$ fois plus de cellules sont exposées au coup suivant~; la
probabilité de la suite complète dans la lignée issue d'une cellule de
départ vaut $g\cdot g^{2}\cdots$ --- pour un $g$ constant par coup, le
risque cumulé devient $N(ut)^{k}
g^{k-1}$ à des facteurs d'ordre près, et l'incidence garde la puissance
$t^{k-1}$ avec un préfacteur $g^{k-1}$ fois plus grand. Si les
expansions croissent elles-mêmes avec le temps (clones à croissance
exponentielle), la courbe se redresse encore, de sorte qu'une pente de
$6$ est produite par moins de sept \omterm{prop:b3:cancer-biology:multistep}{mutations motrices}~: le $k$
d'Armitage--Doll est un majorant du nombre de \omterm{prop:b3:cancer-biology:multistep}{mutations motrices}
limitantes, et les \omterm{def:b3:cancer-biology:hallmarks}{tumeurs} séquencées en montrent de deux à huit.
\end{solution}

\begin{solution}{exo:b3:cancer-biology:11}
(a) Le médicament A seul affronte $\mu_{A}N = 100$ cellules résistantes
préexistantes~: $P_{0} = e^{-100} \approx 0$~; le clone résistant à A
survit et repeuple pendant les dix-huit semaines de A, et lorsque B
commence il forme de nouveau une grande population, de sorte que B
affronte lui aussi $\mu_{B}N' \gg 1$ --- la séquence échoue deux fois.
(b) Ensemble~: une cellule doublement résistante apparaît à $10^{-14}$
par division, $\mu_{A}\mu_{B}N = 10^{-5}$, $P_{0} =
0.99999$. L'association dès le début, quand $N$ est le plus petit, est
la règle parce que c'est le produit de deux petites vitesses qui rend
la résistance improbable.
\end{solution}

\begin{solution}{exo:b3:cancer-biology:12}
Avec $100$ fois plus de cellules et les mêmes $u$ et $k$, le risque
cumulé $N(ut)^{k}$ serait $100$ fois plus élevé --- tout éléphant
devrait mourir jeune d'un \omterm{def:b3:cancer-biology:hallmarks}{cancer}. Ce n'est pas le cas, donc $u$ ou $k$
doivent différer~: vingt copies de \emph{TP53} rendent plus forte la
réponse apoptotique aux lésions et éliminent les cellules mutantes
avant leur expansion (ce qui abaisse le $u$ effectif de cette étape)~;
les grands animaux ont aussi des taux de mutation par cellule plus bas
(meilleure réparation, moins de divisions par cellule et par an, leurs
cellules se renouvelant plus lentement), et peuvent exiger davantage de
coups (une couche de suppresseurs supplémentaire). La suppression des
\omterm{def:b3:cancer-biology:hallmarks}{cancers} a évolué avec la taille du corps.
\end{solution}

\begin{solution}{pb:b3:cancer-biology:1}
\textbf{1.} $10^{12}/10^{3} = 10^{9}$ cellules~; $\log_{2}10^{9} \approx 30$
doublements.
\textbf{2.} $30\times 100 = 3000$ jours, environ $8$ ans.
\textbf{3.} \qty{1}{kg} $\approx 10^{12}$ cellules~: dix doublements de
plus, $1000$ jours, moins de trois ans --- le tiers du temps qu'il a
fallu pour atteindre la détection.
\textbf{4.} Le temps de doublement du début devait valoir environ $60$
jours, ou bien la \omterm{def:b3:cancer-biology:hallmarks}{tumeur} croissait plus vite quand elle était petite et
a ralenti en grossissant (le schéma habituel, l'alimentation et la place
venant à manquer).
\textbf{5.} Un cycle de $2$ jours doublerait la population tous les $2$
jours~; un temps de doublement de $100$ signifie que la croissance
nette vaut $2\,\%$ du taux de naissance~: presque chaque cellule née est
compensée par une qui meurt, ou bien peu de cellules cyclent ---
naissances et morts presque équilibrées, la croissance étant la petite
différence.
\textbf{6.} Les médicaments tuent les cellules en cycle, de sorte
qu'une \omterm{def:b3:cancer-biology:hallmarks}{tumeur} dont une faible fraction est en cycle perd une fraction
moindre par cure~; mais la \omterm{def:b3:cancer-biology:hallmarks}{tumeur} rapide, qui a perdu davantage,
repousse aussi entre les cures. Au total ce sont les \omterm{def:b3:cancer-biology:hallmarks}{tumeurs} rapides
(leucémies, lymphomes, \omterm{def:b3:cancer-biology:hallmarks}{cancer} du testicule) que la chimiothérapie
guérit.
\textbf{7.} Rapport $300$ pour un rapport d'âge $80/30 = 2.67$~: $k - 1 =
\ln 300/\ln 2.67 = 5.8$, $k \approx 7$.
\textbf{8.} $N(ut)^{k} = 10^{8}\times (7\times 10^{-5})^{7} \approx
10^{-21}$~: absurde face à un risque au cours de la vie voisin de $1$
sur $3$. Il manque~: l'expansion clonale qui multiplie les cellules
exposées après chaque coup, des taux de mutation relevés par
l'instabilité, bien plus de divisions (des cellules souches qui cyclent
des décennies) et moins de vraies \omterm{prop:b3:cancer-biology:multistep}{mutations motrices}.
\textbf{9.} Multiplier par $g^{k-1} = 100^{6} = 10^{12}$ donne $10^{-9}$
--- encore minuscule avec sept événements~; avec quatre événements et
des expansions, $10^{8}\times (7\times 10^{-5})^{4}\times 10^{6} \approx 2\times
10^{-3}$, le bon ordre de grandeur. Le tableau vrai est celui de peu de
moteurs et de grandes expansions.
\textbf{10.} $(ut)^{k-1}/(ut)^{k} = 1/(ut) = 1/(10^{-6}\times 40) =
2.5\times 10^{4}$~: l'incidence d'une porteuse à quarante ans vaut des
dizaines de milliers de fois celle des cas sporadiques --- les \omterm{def:b3:cancer-biology:hallmarks}{cancers}
héréditaires frappent jeune.
\textbf{11.} Un événement dix fois plus rapide~: incidence $\times 10$~;
deux~: $\times 100$. Le risque observé, de quinze à vingt-cinq fois,
suggère que la fumée agit sur plus d'une étape.
\textbf{12.} À l'arrêt, $u$ retombe pour les événements qui restent à
venir, de sorte que l'incidence cesse de grimper aussi raide~; mais les
coups déjà accumulés dans les cellules demeurent, de sorte que les
cellules de l'ancien fumeur sont plus près du \omterm{def:b3:cancer-biology:hallmarks}{cancer} que celles d'une
personne n'ayant jamais fumé, et que l'excès de risque persiste, figé
plutôt qu'effacé.
\textbf{13.} $L = \sqrt{2\times 2\times 10^{-9}\times 0.05/0.03} =
\qty{82}{\micro m}$.
\textbf{14.} Entièrement oxygéné jusqu'à un diamètre $2L \approx
\qty{160}{\micro m}$. Un sphéroïde de \qty{2}{mm} a un pourtour vivant
de \qty{82}{\micro m}~: fraction vivante $1 - (918/1000)^{3} = 0.23$.
\textbf{15.} Les cellules de mi-distance sont à \qty{75}{\micro m} d'un
vaisseau, tout juste en deçà de $L$ --- oxygénées en principe, mais les
vaisseaux tumoraux transportent moins de dioxygène et le débit y est
intermittent, de sorte que ces cellules sont hypoxiques. Les cellules
hypoxiques sont deux à trois fois plus résistantes au rayonnement (il
faut du dioxygène pour rendre les dégâts définitifs), et l'hypoxie
déclenche une \omterm{def:b3:cell-cycle-apoptosis:apoptosis}{apoptose} dépendante de \omterm{def:b3:cancer-biology:genes}{p53}, de sorte qu'elle sélectionne
les cellules qui ont perdu \omterm{def:b3:cancer-biology:genes}{p53}.
\textbf{16.} Diviser la densité par deux écarte les vaisseaux d'un
facteur $\sqrt{2}$~: \qty{212}{\micro m} d'écart, cellules de
mi-distance à \qty{106}{\micro m}, au-delà de $L$~: elles meurent ou,
survivant hypoxiques, deviennent plus agressives et plus angiogéniques.
\textbf{17.} $30/2 = 15$ fois plus de glucose. Les \omterm{def:b3:cancer-biology:hallmarks}{tumeurs} concentrent
donc les analogues du glucose, et une scintigraphie au
fluorodésoxyglucose émetteur de positons les fait ressortir.
\textbf{18.} Des \omterm{def:b3:cancer-biology:hallmarks}{tumeurs} affamées se réduisent à la taille limitée par
la diffusion et persistent, dormantes et hypoxiques, ou détournent des
vaisseaux existants~; l'hypoxie sélectionne des clones agressifs. Mais
élaguer la vascularisation chaotique d'une \omterm{def:b3:cancer-biology:hallmarks}{tumeur} abaisse la pression
interstitielle et régularise le débit, de sorte que la chimiothérapie
pénètre mieux --- l'association agit là où ni l'une ni l'autre n'agit
seule.
\textbf{19.} $10^{9}\times 10^{-2n} < 1$~: $n > 4.5$, cinq cures.
\textbf{20.} $21$ jours $= 0.21$ doublement, repousse $\times 1.16$~;
facteur net par cycle $0.0116$~; $\ln 10^{-9}/\ln 0.0116 = 4.65$~:
toujours cinq cycles.
\textbf{21.} $10^{-7}\times 10^{9} = 100$ cellules résistantes
attendues~; $P_{0} = e^{-100} \approx 0$~: la résistance est certaine.
\textbf{22.} $\mu_{1}\mu_{2}N = 10^{-14}\times 10^{9} = 10^{-5}$~;
$P_{0} = 0.99999$. À $10^{12}$ cellules, $10^{-2}$, $P_{0} = 0.99$.
\textbf{23.} $10^{6}$ cellules~: un médicament $\mu N = 0.1$,
$P_{0} = 0.90$~; deux médicaments $10^{-8}$, $P_{0} \approx 1$. Traiter
quand la population est petite, et traiter par des associations.
\textbf{24.} Les lymphocytes~T s'attaquent à beaucoup de néoantigènes à
la fois, de sorte qu'aucune mutation isolée ne leur échappe~; échapper
exige de perdre la présentation antigénique elle-même (CMH,
$\beta_{2}$-microglobuline) ou de réprimer les lymphocytes~T, classe
d'événements différente et plus rare. La même charge mutationnelle
élevée qui rend une \omterm{def:b3:cancer-biology:hallmarks}{tumeur} certaine de résister à un inhibiteur de
kinase lui donne beaucoup de néoantigènes~: les \omterm{def:b3:cancer-biology:hallmarks}{tumeurs} déficientes en
\omterm{prop:b3:dna-repair:fidelity}{réparation des mésappariements}, celles induites par l'ultraviolet et
celles induites par la fumée répondent le mieux.
\textbf{25.} Environ $8$ ans jusqu'à la détection~; de la \omterm{def:b3:cancer-biology:hallmarks}{tumeur} vivante
sur \qty{82}{\micro m} autour d'un vaisseau~; $P_{0} = 0.99999$ qu'une
association de deux médicaments n'affronte aucune cellule résistante
dans une \omterm{def:b3:cancer-biology:hallmarks}{tumeur} de \qty{1}{cm^3}.
\end{solution}
